% Shared definitions for study-guide.tex and exercises.tex
% (Week 9).
% Same colors and box look as script.tex; no lecturer cues.

\definecolor{navy}{RGB}{24,48,75}
\definecolor{teal}{RGB}{0,120,125}
\definecolor{gold}{RGB}{190,130,20}

\addtokomafont{disposition}{\color{navy}}
\setkomafont{section}{\Large\bfseries\color{navy}}

\newcommand{\Om}{\Omega}
% Expectation, variance and standard deviation, as in Ross: E[X], Var(X), SD(X)
\newcommand{\E}{\operatorname{E}}
\newcommand{\Var}{\operatorname{Var}}
\newcommand{\SD}{\operatorname{SD}}
\newcommand{\myemail}{tiborsandor2001@gmail.com}

% Exercise statement (identical in both versions)
\newtcolorbox{task}[1]{enhanced,breakable,colback=navy!4,colframe=navy,
	fonttitle=\bfseries,title={#1}}
% Finance example
\newtcolorbox{fin}[1]{enhanced,breakable,colback=teal!5,colframe=teal,
	fonttitle=\bfseries,title={#1}}
% Key formulas at the start of each part
\newtcolorbox{key}[1]{enhanced,breakable,colback=white,colframe=navy!70,
	fonttitle=\bfseries,title={#1}}
% Final answer (solutions version only)
\newtcolorbox{answer}{enhanced,colback=teal!6,colframe=teal,boxrule=.5pt,
	left=4pt,right=4pt,top=2pt,bottom=2pt}
% Remarks: common mistakes, alternative methods
\newtcolorbox{remark}{enhanced,breakable,colback=gold!8,colframe=gold,
	boxrule=.4pt,left=4pt,right=4pt,top=2pt,bottom=2pt,fontupper=\small}

% Labels inside solutions
\newcommand{\soltitle}{\par\textbf{\color{teal}Solution.}\par\nobreak}
\newcommand{\idea}{\textbf{Idea.}\ }
\newcommand{\step}[1]{\par\textbf{#1}\ }

% Small Venn diagram of two events E, F in Omega; #1 = shading commands
\newcommand{\venn}[1]{%
	\begin{tikzpicture}[scale=.5,baseline=(O.base)]
		\node (O) at (.5,0) {\phantom{x}};
		#1
		\draw[thick] (-1.5,-1.15) rectangle (2.5,1.15);
		\draw[thick,navy] (0,0) circle (.85);
		\draw[thick,navy] (1,0) circle (.85);
		\node at (-.35,0) {\scriptsize $E$};
		\node at (1.35,0) {\scriptsize $F$};
		\node at (2.15,.8) {\scriptsize $\Omega$};
	\end{tikzpicture}}
\newcommand{\shadeall}[1]{\fill[#1] (-1.5,-1.15) rectangle (2.5,1.15);}

% Visual proof of Bayes' formula: Omega = gray box, the hypothesis F = blue
% diamond, the observation E = light blue circle, their overlap FE = red.
% Mini shapes (pics): bvOm, bvF, bvE, bvFE; bvfrac={top}{bottom} is a
% fraction of two mini shapes. Each row is one tikzpicture, so the HTML
% conversion sees no math around TikZ. Optional argument: scale.
\definecolor{bvblue}{RGB}{95,107,216}
\definecolor{bvlight}{RGB}{214,228,250}
\definecolor{bvrose}{RGB}{248,122,134}
\definecolor{bvgray}{RGB}{231,234,237}
\tikzset{
  bvline/.style={draw=black!80,line width=.5pt},
  bvclip/.style={draw=black!80,line width=1pt},
  bvstrike/.style={draw=red!85!black,line width=1pt},
  pics/bvOm/.style={code={\filldraw[bvline,fill=bvgray,rounded corners=2pt]
    (-.3,-.16) rectangle (.3,.16);}},
  pics/bvF/.style={code={\filldraw[bvline,fill=bvblue]
    (0,.21)--(.21,0)--(0,-.21)--(-.21,0)--cycle;}},
  pics/bvE/.style={code={\filldraw[bvline,fill=bvlight] (0,0) circle (.2);}},
  pics/bvFE/.style={code={%
    \begin{scope}[scale=.25,shift={(-.12,0)}]
      \clip (-.6,1.25)--(.65,0)--(-.6,-1.25)--(-1.85,0)--cycle;
      \filldraw[bvclip,fill=bvrose] (.8,0) circle (1.2);
    \end{scope}
    \begin{scope}[scale=.25,shift={(-.12,0)}]
      \clip (.8,0) circle (1.2);
      \draw[bvclip] (-.6,1.25)--(.65,0)--(-.6,-1.25)--(-1.85,0)--cycle;
    \end{scope}}},
  pics/bvfrac/.style 2 args={code={\pic at (0,.3) {#1};
    \draw[line width=.5pt] (-.36,0)--(.36,0); \pic at (0,-.3) {#2};}},
  pics/bvx/.style={code={\draw[bvstrike] (-.3,-.21)--(.3,.21);}},
}
% Row 1: Omega, F, E and the overlap FE
\newcommand{\bayesvenn}[1][1]{%
  \begin{tikzpicture}[scale=#1]
    \filldraw[bvline,line width=.8pt,fill=bvgray,rounded corners=6pt]
      (-2.4,-1.6) rectangle (2.4,1.6);
    \filldraw[bvline,line width=.8pt,fill=bvlight] (.8,0) circle (1.2);
    \filldraw[bvline,line width=.8pt,fill=bvblue]
      (-.6,1.25)--(.65,0)--(-.6,-1.25)--(-1.85,0)--cycle;
    \begin{scope}
      \clip (-.6,1.25)--(.65,0)--(-.6,-1.25)--(-1.85,0)--cycle;
      \fill[bvrose] (.8,0) circle (1.2);
    \end{scope}
    \draw[bvline,line width=.8pt] (.8,0) circle (1.2)
      (-.6,1.25)--(.65,0)--(-.6,-1.25)--(-1.85,0)--cycle;
    \node[white,font=\large] at (-.85,0) {$F$};
    \node[font=\large] at (1.3,0) {$E$};
    \node[font=\footnotesize] at (.13,0) {$FE$};
    \node at (2.05,1.25) {$\Omega$};
  \end{tikzpicture}}
% Row 2: probability = share of a shape
\newcommand{\bayesshares}[1][1]{%
  \begin{tikzpicture}[scale=#1]
    \node[left] at (0,0) {$P(F)=$};
    \pic at (.42,0) {bvfrac={bvF}{bvOm}};
    \node[left] at (2.55,0) {$P(E)=$};
    \pic at (2.97,0) {bvfrac={bvE}{bvOm}};
    \node[left] at (5.55,0) {$P(F\mid E)=$};
    \pic at (5.97,0) {bvfrac={bvFE}{bvE}};
    \node[left] at (8.55,0) {$P(E\mid F)=$};
    \pic at (8.97,0) {bvfrac={bvFE}{bvF}};
  \end{tikzpicture}}
% Row 3: plug the shares into Bayes' formula; F and Omega cancel
\newcommand{\bayesproof}[1][1]{%
  \begin{tikzpicture}[scale=#1]
    \node[left] at (0,0)
      {$\displaystyle P(F\mid E)=\frac{P(E\mid F)\,P(F)}{P(E)}=$};
    \pic at (.8,.72) {bvfrac={bvFE}{bvF}};
    \node at (1.4,.72) {$\cdot$};
    \pic at (2,.72) {bvfrac={bvF}{bvOm}};
    \draw[line width=.6pt] (.3,0)--(2.5,0);
    \pic at (1.4,-.72) {bvfrac={bvE}{bvOm}};
    \pic at (.8,.42) {bvx};
    \pic at (2,1.02) {bvx};
    \pic at (2,.42) {bvx};
    \pic at (1.4,-1.02) {bvx};
    \node at (2.9,0) {$=$};
    \pic at (3.55,0) {bvfrac={bvFE}{bvE}};
    \node[right] at (3.95,0) {$=P(F\mid E)$};
    \node[right,font=\footnotesize] at (3.95,-.45) {(the definition)};
  \end{tikzpicture}}

\title{Mathematics for Finance}
\author{Tibor SÁNDOR\\[2pt]
	\normalsize Corvinus University of Budapest\\
	\normalsize\href{mailto:\myemail}{\texttt{\myemail}}}

% The exercises (one file each in exercises/) and the commands that place them
% Exercises of Week 9, shared by beamer.tex, the study guide, the script and
% exercises.tex. Every exercise lives in a file of its own,
% exercises/NN-slug.tex, numbered 9.NN (see exercises.lua for the docstring).
% The body of an exercise file:
%
%   \begin{problem} ... \end{problem}      the statement, used everywhere;
%                                          inside it, \begin{context} ...
%                                          \end{context} is background that
%                                          the slides leave out
%   \begin{slideproblem} ... \end{slideproblem}
%                                          (optional) a shorter statement
%                                          for the slides
%   \begin{slides} ... \end{slides}        one frame of the deck: \exstatement
%                                          places the statement; repeat the
%                                          environment for more frames
%   \begin{solution} ... \end{solution}    the detailed solution (study
%                                          guide, exercise PDFs, website)
%
% Commands for the documents (the argument is the slug: 01-license-plates.tex
% has the slug license-plates; an exercise of an earlier week included below
% is referred to the same way):
%
%   \exref{slug}                 the number, e.g. 9.1, for "Exercise~\exref{...}"
%   \extitle{slug}               the title from the docstring
%   \exerciseframe[opts]{slug}   beamer: the frames of the exercise
%   \studyexercise{slug}         study guide: heading, statement, solution
%   \scriptexercise{slug}        script: the statement box, as on the slide
%   \allexercises                exercises.tex: every exercise, in order; on
%                                the website statement and solution with
%                                markers where the pipeline splits the HTML
%                                (\collectionexercise), in the PDF plain:
%                                each statement, followed by its solution if
%                                \ifsolutions (exercises-solutions.tex)

\def\exweek{9}

\directlua{exercises = dofile("exercises/exercises.lua")
  exercises.scan("exercises", "\exweek")}

% Exercises of earlier weeks used in this week's documents (by slug, numbered
% as in their own week, e.g. 8.11; not listed among this week's exercises).
% scripts/build-content.mjs reads these lines to copy the files.
\directlua{exercises.include("week-08")}

\ExplSyntaxOn

\msg_new:nnn { exercises } { unknown } { Unknown~exercise~`#1'. }

\cs_new_protected:Npn \__ex_missing:n #1
  { \msg_error:nnn { exercises } { unknown } {#1} ?? }
\cs_new:Npn \exref #1
  {
    \cs_if_exist:cTF { exnum@#1 }
      { \use:c { exnum@#1 } } { \__ex_missing:n {#1} }
  }
\cs_new:Npn \extitle #1
  {
    \cs_if_exist:cTF { extitle@#1 }
      { \use:c { extitle@#1 } } { \__ex_missing:n {#1} }
  }

% ---------------------------------------------------------------------------
% Loading an exercise: its parts are stored, nothing is typeset. The slides
% typeset the stored parts; the other documents read the part they need from
% the file again (\__ex_input_part:nn), as stored tables do not survive lwarp.

\tl_new:N \g__ex_current_tl
\tl_new:N \g__ex_problem_tl
\tl_new:N \g__ex_slideproblem_tl
\tl_new:N \g__ex_solution_tl
\seq_new:N \g__ex_slides_seq
\box_new:N \l__ex_scratch_box

\cs_new_protected:Npn \__ex_load:n #1
  {
    \cs_if_exist:cF { exfile@#1 } { \__ex_missing:n {#1} }
    \tl_gset:Nn \g__ex_current_tl {#1}
    \tl_gclear:N \g__ex_problem_tl
    \tl_gclear:N \g__ex_slideproblem_tl
    \tl_gclear:N \g__ex_solution_tl
    \seq_gclear:N \g__ex_slides_seq
    \vbox_set:Nn \l__ex_scratch_box
      {
        \DeclareDocumentEnvironment { problem } { +b }
          { \tl_gset:Nn \g__ex_problem_tl {##1} } { }
        \DeclareDocumentEnvironment { slideproblem } { +b }
          { \tl_gset:Nn \g__ex_slideproblem_tl {##1} } { }
        \DeclareDocumentEnvironment { slides } { +b }
          { \seq_gput_right:Nn \g__ex_slides_seq {##1} } { }
        \DeclareDocumentEnvironment { solution } { +b }
          { \tl_gset:Nn \g__ex_solution_tl {##1} } { }
        \exp_args:Nv \input { exfile@#1 }
      }
  }

% Typesets one part (problem, slideproblem, solution) of an exercise straight
% from its file; the other parts are skipped.
\cs_new_protected:Npn \__ex_input_part:nn #1#2
  {
    \group_begin:
      \clist_map_inline:nn { problem, slideproblem, slides, solution }
        {
          \str_if_eq:nnTF {##1} {#2}
            { \DeclareDocumentEnvironment {##1} { } { } { } }
            { \DeclareDocumentEnvironment {##1} { +b } { } { } }
        }
      \exp_args:Nv \input { exfile@#1 }
    \group_end:
  }

% Exercise 9.1, or Exercise 9.2 · Finance
\cs_new:Npn \__ex_heading:
  {
    Exercise~\exref { \g__ex_current_tl }
    \cs_if_exist:cT { exfinance@ \g__ex_current_tl } { ~·~Finance }
  }
\prg_new_conditional:Npnn \__ex_if_finance: { T, F, TF }
  {
    \cs_if_exist:cTF { exfinance@ \g__ex_current_tl }
      { \prg_return_true: } { \prg_return_false: }
  }

% ---------------------------------------------------------------------------
% Beamer

% The finance context is printed in the study guide, the exercise PDFs and on
% the website; the
% slides and the script (which shows what the slide shows) leave it out.
\bool_new:N \l__ex_context_bool
\bool_set_true:N \l__ex_context_bool
\DeclareDocumentEnvironment { context } { }
  {
    \bool_if:NTF \l__ex_context_bool
      { \par \textit { Finance~context. } ~ }
      { \__ex_gobble_context:w }
  }
  { \bool_if:NT \l__ex_context_bool { \par } }
\cs_new_protected:Npn \__ex_gobble_context:w #1 \end #2
  { \str_if_eq:nnTF {#2} { context } { \end { context } } { \__ex_gobble_context:w } }

\IfClassLoadedTF { beamer }
  {
    \bool_set_false:N \l__ex_context_bool

    \cs_new:Npn \__ex_block: { \__ex_if_finance:TF { exampleblock } { block } }
    \NewDocumentCommand \exstatement { }
      {
        \begin { \__ex_block: } { \__ex_heading: }
          \tl_if_empty:NTF \g__ex_slideproblem_tl
            { \g__ex_problem_tl } { \g__ex_slideproblem_tl }
        \end { \__ex_block: }
      }

    % The steps (\pause, \uncover, ...) of an exercise are not shown one by
    % one: the first frame gives two pages, its first slide (the statement,
    % no solution) and its last state (the full solution); a further frame
    % (cont.) only its last state. Slide 99 is past every step, so it shows
    % everything (a frame has only the slides it uses, so <1,99> is two).
    \tl_new:N \l__ex_frames_tl
    \NewDocumentCommand \exerciseframe { O{t} m }
      {
        \__ex_load:n {#2}
        \tl_clear:N \l__ex_frames_tl
        \seq_map_indexed_inline:Nn \g__ex_slides_seq
          {
            \int_compare:nNnTF {##1} = { 1 }
              { \tl_put_right:Nn \l__ex_frames_tl { \begin{frame}<1,99>[#1] } }
              { \tl_put_right:Nn \l__ex_frames_tl { \begin{frame}<99>[#1] } }
            \tl_put_right:Ne \l__ex_frames_tl
              {
                {
                  \exp_not:N \extitle {#2}
                  \int_compare:nNnT {##1} > { 1 } { \c_space_tl (cont.) }
                }
              }
            \tl_put_right:Nn \l__ex_frames_tl { ##2 \end{frame} }
          }
        \tl_use:N \l__ex_frames_tl
      }
  }
  {
    \cs_new:Npn \__ex_box: { \__ex_if_finance:TF { fin } { task } }
    \cs_new_protected:Npn \__ex_statement_box:n #1
      {
        \begin { \__ex_box: } { \__ex_heading: }
          \__ex_input_part:nn { \g__ex_current_tl } {#1}
        \end { \__ex_box: }
      }

    \NewDocumentCommand \studyexercise { m }
      {
        \__ex_load:n {#1}
        \subsection* { Exercise~\exref{#1}~---~\extitle{#1} }
        \addcontentsline { toc } { subsection }
          { Exercise~\exref{#1}~---~\extitle{#1} }
        \__ex_statement_box:n { problem }
        \soltitle \__ex_input_part:nn {#1} { solution }
      }

    \NewDocumentCommand \scriptexercise { m }
      {
        \__ex_load:n {#1}
        \group_begin:
          \bool_set_false:N \l__ex_context_bool
          \tl_if_empty:NTF \g__ex_slideproblem_tl
            { \__ex_statement_box:n { problem } }
            { \__ex_statement_box:n { slideproblem } }
        \group_end:
      }

    % The website (lwarp): statement and solution of each exercise, between
    % markers where the pipeline splits the HTML.
    \NewDocumentCommand \collectionexercise { m }
      {
        \__ex_load:n {#1}
        \par [[exercise:#1:problem]] \par
        \__ex_input_part:nn {#1} { problem }
        \par [[exercise:#1:solution]] \par
        \__ex_input_part:nn {#1} { solution }
        \par [[exercise:#1:end]] \par
      }
    % The PDF: plain, each statement, followed by its solution in the
    % version with solutions (\ifsolutions). The statement is set off like a
    % quote in Markdown, by a bar on the left (quote=<color>, also used by
    % exercises.tex) in the color of its box in the study guide: navy, or teal
    % for a finance exercise.
    \tcbset
      {
        quote/.style =
          {
            enhanced, breakable, frame~hidden, interior~hidden,
            sharp~corners, boxrule = 0pt, borderline~west = {2pt}{0pt}{#1},
            left = 8pt, right = 0pt, top = 1pt, bottom = 1pt,
            before~skip = 6pt, after~skip = 6pt,
            % the document's paragraph spacing (parskip=half)
            before~upper = { \setlength { \parskip } { .5\baselineskip } }
          }
      }
    \NewDocumentCommand \plainexercise { m }
      {
        \__ex_load:n {#1}
        \subsection* { Exercise~\exref{#1}~---~\extitle{#1} }
        \addcontentsline { toc } { subsection }
          { Exercise~\exref{#1}~---~\extitle{#1} }
        \__ex_if_finance:TF
          { \begin{tcolorbox}[quote=teal] }
          { \begin{tcolorbox}[quote=navy] }
          \__ex_input_part:nn {#1} { problem }
        \end{tcolorbox}
        \legacy_if:nT { solutions }
          {
            \soltitle
            \__ex_input_part:nn {#1} { solution }
          }
      }
    \NewDocumentCommand \allexercises { }
      {
        \IfPackageLoadedTF { lwarp }
          { \directlua { exercises.each("collectionexercise") } }
          { \directlua { exercises.each("plainexercise") } }
      }
  }

\ExplSyntaxOff


% PDF properties and the date on the title page (\documentinfo)
% The properties of every PDF of Week 9: author (with email address), title,
% subject, keywords and the date the document last changed, in the PDF
% (hyperref) and on its title page (\date), and the header and footer of the
% printed documents (below). Each document loads this file after hyperref and
% then calls \documentinfo{<what it is>}, e.g. \documentinfo{Slides}.
%
% The date: the website build (scripts/build-content.mjs) puts
% \def\lastmodified{YYYY-MM-DD} in front of the document, the date of the last
% commit that changed its sources. Compiled locally, it is today.

\newcommand{\courseinfo}{Mathematics for Finance}
\newcommand{\weeklabel}{Week 9}
\newcommand{\weekinfo}{\weeklabel: Probability, conditional probability, independence, random variables}
\newcommand{\authorname}{Tibor SÁNDOR}
\providecommand{\myemail}{tiborsandor2001@gmail.com}

\makeatletter
\ifdefined\lastmodified\else
	\edef\lastmodified{\the\year-\two@digits{\the\month}-\two@digits{\the\day}}
\fi
\def\di@split#1-#2-#3\@nil{%
	\def\di@year{#1}\def\di@month{#2}\def\di@day{#3}}
\expandafter\di@split\lastmodified\@nil
% November 11, 2026
\edef\lastmodifiedtext{%
	\ifcase\numexpr\di@month\relax\or January\or February\or March\or April\or
		May\or June\or July\or August\or September\or October\or November\or
		December\fi\space\number\di@day, \di@year}
% D:20261111000000, the date format of PDF metadata
\edef\di@pdfdate{D:\di@year\di@month\di@day000000}

\newcommand{\documentinfo}[1]{%
	\hypersetup{
		pdftitle={\courseinfo, \weekinfo\space -- #1},
		pdfauthor={\authorname\space<\myemail>},
		pdfsubject={\courseinfo, \weekinfo},
		pdfkeywords={probability, combinatorics, mathematics for finance,
				Corvinus University of Budapest},
		pdfcreationdate={\di@pdfdate},
		pdfmoddate={\di@pdfdate},
		pdflang={en-US},
	}%
	\date{Last updated: \lastmodifiedtext}%
	\gdef\documentkind{#1}%
}

% Header and footer of the printed documents (KOMA-Script, scrlayer-scrpage;
% the slides have beamer's footline, the website neither): the course, the
% week and the document on the left of the header (the week's topic is on the
% title page), the university on the right; in the footer the author and the
% page as <current>/<total>. The title page gets only the footer.
\newcommand{\documentkind}{}
\IfClassLoadedTF{beamer}{}{\IfPackageLoadedTF{lwarp}{}{%
	\usepackage{lastpage}%
	\usepackage[headsepline=1.5pt,footsepline=1.5pt,plainfootsepline]
		{scrlayer-scrpage}%
	\clearpairofpagestyles
	\setkomafont{pageheadfoot}{\small\sffamily\color{navy}}%
	\setkomafont{pagenumber}{}%
	\setkomafont{headsepline}{\color{navy!35}}%
	\setkomafont{footsepline}{\color{navy!35}}%
	\ihead{\courseinfo\ \textperiodcentered\ \weeklabel\ \textperiodcentered\ 
		\documentkind}%
	\ohead{Corvinus University of Budapest}%
	\ifoot*{\authorname\ \textperiodcentered\ \myemail}%
	\ofoot*{\pagemark/\pageref*{LastPage}}%
	\pagestyle{scrheadings}%
	% Extra stretch for breaking lines with long inline formulas.
	\setlength{\emergencystretch}{1.5em}%
}}
\makeatother

